\documentclass{szb-solution}

\area{Integrálszámítás}
\title{Integrálszámítás IV}
\subject{Matematika G1}
\subjectCode{BMETE94BG01}
\date{Utoljára frissítve: \today}
\docno{13}

\begin{document}
\maketitle

\subsection{Paraméteres határozott integrál minimuma}

Legyen $f(x)=x^4-2x^2=x^2(x^2-2)$. A függvény negatív pontosan a
$(-\sqrt2,\sqrt2)$ intervallumon, a végpontokban nulla, azon kívül pozitív.
Ezért az
$$
  I(a,b)=\int_a^b f(x)\dd x,
  \qquad a<b,
$$
integrál akkor a legkisebb, ha a teljes negatív részt integráljuk, pozitív
részt pedig nem veszünk hozzá. Tehát
$$
  \boxed{a=-\sqrt2,\qquad b=\sqrt2}
  \text.
$$
Az integrandus páros, így a minimális érték
\begin{align*}
  I_{\min}
  &=2\int_0^{\sqrt2}(x^4-2x^2)\dd x\\
  &=2\left[\frac{x^5}{5}-\frac{2x^3}{3}\right]_0^{\sqrt2}
  =\boxed{-\frac{16\sqrt2}{15}}
  \text.
\end{align*}

\subsection{Improprius integrálok}

\begin{enumerate}[a)]
  \item A felső határ végtelen:
        $$
          \int_0^\infty\frac{\dd x}{1+x^2}
          =\lim_{R\to\infty}[\arctan x]_0^R
          =\boxed{\frac\pi2}
          \text.
        $$
        Az integrál konvergens.

  \item A függvény a bal végpontban nem korlátos, ezért
        $$
          \int_0^1\frac{\dd x}{\sqrt x}
          =\lim_{\varepsilon\to0^+}[2\sqrt x]_{\varepsilon}^{1}
          =\boxed{2}
          \text.
        $$
        Az integrál konvergens.

  \item A bal végpontban pólus van:
        \begin{align*}
          \int_{-2}^0\frac{\dd x}{x+2}
          &=\lim_{\varepsilon\to0^+}
            [\ln(x+2)]_{-2+\varepsilon}^{0}\\
          &=\lim_{\varepsilon\to0^+}(\ln2-\ln\varepsilon)=+\infty
          \text.
        \end{align*}
        Ezért az integrál \boxed{\text{divergens}}.

  \item Legyen $t=\sqrt{x-1}$, vagyis $x=t^2+1$ és
        $\dd x=2t\dd t$. Ekkor
        $$
          \int_1^\infty\frac{\dd x}{x\sqrt{x-1}}
          =2\int_0^\infty\frac{\dd t}{1+t^2}
          =\boxed{\pi}
          \text.
        $$

  \item Egy primitív függvény
        $$
          \int e^{-x}\cos x\dd x
          =\frac{e^{-x}}2(\sin x-\cos x)+C
          \text.
        $$
        Mivel az exponenciális tényező a végtelenben nullához tart,
        $$
          \int_0^\infty e^{-x}\cos x\dd x
          =\boxed{\frac12}
          \text.
        $$

  \item A bal végpontban $x\ln x\to0$, ezért
        $$
          \int_0^1\ln x\dd x
          =\lim_{\varepsilon\to0^+}[x\ln x-x]_{\varepsilon}^{1}
          =\boxed{-1}
          \text.
        $$
\end{enumerate}

\subsection{A parabolaív hossza}

Az $f(x)=x^2$ függvény deriváltja $f'(x)=2x$. Az ívhossz
$$
  L=\int_0^2\sqrt{1+4x^2}\dd x
  \text.
$$
Az $u=2x$ helyettesítéssel
$$
  L=\frac12\int_0^4\sqrt{1+u^2}\dd u
  \text.
$$
Felhasználva, hogy
$$
  \int\sqrt{1+u^2}\dd u
  =\frac12\left(u\sqrt{1+u^2}+\operatorname{arsinh}u\right)+C,
$$
kapjuk
$$
  \boxed{L=\sqrt{17}+\frac14\operatorname{arsinh}4
  =\sqrt{17}+\frac14\ln(4+\sqrt{17})}
  \text.
$$

\subsection{A ciklois egy ívének hossza}

A paraméterezés deriváltjai
$$
  \dot x(t)=a(1-\cos t),
  \qquad
  \dot y(t)=a\sin t
  \text.
$$
Ezért a sebesség nagysága
\begin{align*}
  \sqrt{\dot x^2+\dot y^2}
  &=a\sqrt{(1-\cos t)^2+\sin^2t}\\
  &=a\sqrt{2-2\cos t}
   =2a\left|\sin\frac t2\right|
  \text.
\end{align*}
A $[0,2\pi]$ intervallumon $\sin(t/2)\geq0$, így
$$
  L=2a\int_0^{2\pi}\sin\frac t2\dd t
  =\boxed{8a}
  \text.
$$

\subsection{A csonkakúp térfogata}

Vegyük fel az $x$ tengelyt a csonkakúp tengelye mentén. A sugár a
$[0,6]$ intervallumon lineárisan csökken $5$-ről $2$-re:
$$
  r(x)=5-\frac{x}{2}
  \text.
$$
Korongmódszerrel
\begin{align*}
  V
  &=\pi\int_0^6\left(5-\frac x2\right)^2\dd x\\
  &=\pi\left[25x-\frac52x^2+\frac{x^3}{12}\right]_0^6
  =\boxed{78\pi}
  \text.
\end{align*}
Ez megegyezik a
$V=\pi h(R^2+Rr+r^2)/3$ képletbe való közvetlen behelyettesítés
eredményével.

\subsection{A gömb térfogatképletének levezetése}

Az $R$ sugarú gömböt az
$$
  y=\sqrt{R^2-x^2},
  \qquad -R\leq x\leq R,
$$
felső félkör $x$ tengely körüli megforgatásával kapjuk. A korongmódszer
alapján
\begin{align*}
  V
  &=\pi\int_{-R}^{R}y^2\dd x
    =\pi\int_{-R}^{R}(R^2-x^2)\dd x\\
  &=\pi\left[R^2x-\frac{x^3}{3}\right]_{-R}^{R}
    =\boxed{\frac43\pi R^3}
  \text.
\end{align*}

\subsection{Síktartomány súlypontja}

Az $f(x)=x^3$, az $x$ tengely és az $x=1$ egyenes által határolt tartomány
$0\leq x\leq1$ között helyezkedik el. Területe
$$
  T=\int_0^1x^3\dd x=\frac14
  \text.
$$
A súlypont vízszintes koordinátája
$$
  S_x=\frac{\int_0^1x\,x^3\dd x}{T}
     =\frac{1/5}{1/4}=\frac45
  \text,
$$
a függőleges koordinátája pedig
$$
  S_y=\frac{\frac12\int_0^1(x^3)^2\dd x}{T}
     =\frac{1/14}{1/4}=\frac27
  \text.
$$
Tehát
$$
  \boxed{S\left(\frac45,\frac27\right)}
  \text.
$$

\end{document}
